\section{Registration responds to atypicality, not to lead}\label{sec:registration}\label{sec:reg_rates}\label{sec:representation}\label{sec:amend}

The first step sets up the contrast with every step after it. The
companion paper shows, with deciles cut within Nice class and filing
year, that the registration rate rises with atypicality from 56.0\% to
59.1\% and turns down to 58.3\% in the most atypical tenth, while across
deciles of lead it moves by less than a point \citep{SilverCorpus2026}.
That flatness has a plain reading: lead is computed from filings made in
the five years \emph{after} an application, which had not happened when
the application was examined, so no one in prosecution could have
responded to it. How unusual a description is, by contrast, is exactly
what an examiner steeped in the class's usual language notices.
Registration is also mostly a measure of the applicant's own
follow-through rather than an examiner's verdict, since nine in ten
failures are the applicant's own inaction.

Two facts about this step carry into the rest of the paper. Who drafts
the description matters: represented applications (73\% of filings)
register at 61.9\% against 45.6\% for self-filed ones, and self-filers
write shorter descriptions that sit closer to their class's usual
language, so counsel of record is held in every five-year-proof
regression and the estimates are reported separately by drafter. And
applicants who try again learn to conform: among abandoned applications
later registered by the same owner in the same class, three in four
carry a rewritten description, moved toward the class's typical
language from both tails. Every estimate on registered marks is
therefore an estimate on partly conformed language.
